\section{Introduction}\label{sec:intro}

Hybrid post-quantum/traditional (PQ/T) key exchange pairs a post-quantum KEM
such as ML-KEM~\cite{FIPS203} with a Diffie--Hellman group, so that the
result stays secure as long as either one does. X-Wing~\cite{XWing,XWingDraft}
fixes one such pairing, ML-KEM-768 with X25519, and proves it IND-CCA secure;
generic KEM combiners~\cite{GHP18,BBFGS19} cover other pairings; and composite
signatures~\cite{Composite} pair ML-DSA~\cite{FIPS204} with EdDSA~\cite{RFC8032}
in the same spirit. A protocol built from these primitives inherits their
guarantees only through a composition argument, and that argument is usually
made on paper.

This paper reports a machine-checked computational analysis, in
CryptoVerif~\cite{Bla08,CryptoVerif}, of HYBRID-AKE: a small SIGMA-style
handshake~\cite{Kra03} that derives its key from an ephemeral hybrid KEM and
authenticates with composite signatures over the transcript. We mechanize
both layers. The protocol proof assumes one IND-CCA2 KEM and one UF-CMA
signature; separate proofs show that each hybrid primitive meets that
assumption if either of its components does. The protocol proof and the
post-quantum legs use only CryptoVerif's post-quantum-sound
library~\cite{BJ24}.

The protocol is not the point. It is an engineering design, close to TLS~1.3
and SIGMA, and it is not standardised. What we report is what mechanization,
negative controls and an independent review caught, at three levels:
\begin{enumerate}
  \item \emph{The proof sketch.} A pen-and-paper sketch preceded the
  mechanization. Two of its four results were stated without an assumption
  they need. The composite signature is unforgeable when either component is
  only if its pre-hash is also collision resistant. X-Wing is secure relying
  on X25519 only if ML-KEM is also ciphertext second-preimage resistant
  (C2PRI). Both conditions are stated in the primary sources, the
  composite-signature draft~\cite{Composite} and the X-Wing paper's
  classical theorem~\cite{XWing}; the sketch had dropped them.
  \item \emph{Our own write-up.} An independent review of the first version
  of the mechanization found claims that the models did not support: a
  tighter bound than proved, non-standard assumptions not flagged as such,
  and an adversary weaker than described.
  \item \emph{The implementation.} Comparing the models' idealised X25519
  with the code showed that our X-Wing implementation raised an error on
  low-order X25519 points, where the X-Wing specification's algorithm
  returns a shared secret. The same divergence had been reported for another
  library~\cite{Noble26}.
\end{enumerate}

\paragraph{Contributions.}
\begin{itemize}
  \item Machine-checked proofs (9 models, 17 queries, CryptoVerif~2.13, under
  20~seconds in total) of: injective agreement in both directions;
  real-or-random secrecy of every fresh session key, with forward secrecy,
  key-compromise impersonation resistance and key reveal; cross-role key
  uniqueness; the same results under static ephemeral-key reveal; IND-CCA2 of
  X-Wing relying on either leg; IND-CCA2 of an ML-KEM-1024/X25519
  transcript-hashing combiner relying on either leg, including a
  random-oracle proof from standard assumptions; and EUF-CMA of the composite
  signature relying on either component (\cref{sec:model,sec:mech}).
  \item Sixteen mutation tests, each deleting one design element or
  assumption, with controls (\cref{sec:mutations}). They show that the KEM
  ciphertext and the responder's key must be in the signed transcript, that
  every freshness guard in the protocol model has an effect, that X-Wing
  needs C2PRI where the transcript-hashing combiner does not, and that the
  role-separating context strings are redundant given the encodings.
  \item A catalogue of the findings above, with their causes and fixes, and
  the lessons we draw from them (\cref{sec:caught}).
  \item For the implementation finding, a fix and a complete
  characterisation of the X25519 inputs that produce an all-zero output,
  checked against an independent RFC~7748 implementation and by counting
  roots over $\Fp$ (\cref{sec:lowcaught,app:loworder}).
\end{itemize}
All models, proof scripts, mutation tests, implementation tests and the
sources of this paper are public~\cite{Artifact}.

We are not aware of earlier machine-checked proofs of X-Wing or of composite
ML-DSA/EdDSA signatures. The closest mechanized work covers generic KEM
combiners and KEM-based hybrid key exchanges in Squirrel~\cite{BDDJK26}, and
the hybrid TLS~1.3 and SSH key exchanges in CryptoVerif~\cite{BJ24}
(\cref{sec:related}).

\paragraph{How this work was done.}
The models, proof scripts, tests and a first draft of this paper were
written with an AI assistant (Claude, by Anthropic) at the author's
direction. The independent review of \cref{sec:reviewcaught} was carried out
by a separate instance of the same assistant that had not seen the
development. We state this because it bears on how the results should be
weighed. The machine-checked statements rest on CryptoVerif's soundness and
on the models matching the specification. The review, by contrast, is not a
substitute for expert peer review (\cref{sec:limits}).

\paragraph{Outline.}
\Cref{sec:background} describes the protocol, its primitives and
CryptoVerif. \Cref{sec:model} states the security model and the results,
marking which parts are machine-checked. \Cref{sec:mech} explains the
mechanization, \cref{sec:mutations} the mutation tests, and
\cref{sec:caught} what was caught. \Cref{sec:limits} lists the limitations,
\cref{sec:related} discusses related work, and \cref{sec:conclusion}
concludes. The appendices give model excerpts, the low-order analysis and
reproduction instructions.
